\section{Selection and continuation on the initial time interval}
\label{sec:early}

The cold estimates hold without an a priori bound on the hot part.
We use them first to select a hot seed whose evolution touches a prescribed
cap, and then to continue every selected evolution to a fixed time $t_1$.
The second step must be uniform as the cap horizon $T$ tends to zero.
The bath remains weak inside the hole at very small times; away from time
zero, the thinness of the jet tubes gives the additional smallness we need.
All constants in this section may depend on the fixed jet parameters.
They are independent of $N$, the collision cutoff, and the seed parameter,
unless explicitly stated otherwise.

\subsection{The bath and the initial source}

\begin{lemma}[Bath barrier]\label{early:bath}
There are $t_1\in(0,1/32]$, $\beta'>4$, and constants $C,c>0$,
independent of $\delta$, $R_\mu$, and $s_0$, such that every finite evolution
of \eqref{con:colors} satisfies, on its interval of existence,
\begin{equation}\label{early:bath-bound}
 0\le u(t,x,v)\le C e^{-\beta'|v|^2}
       \exp\!\left(-\frac{cD_0(x)^2}{t}\right),
 \qquad 0<t\le t_1,
 \qquad D_0(x)=\dist(x,D_{\rm hole}^c).
\end{equation}
For sufficiently small $s_0$ depending on $\delta$, one also has
\begin{equation}\label{early:hole-distance}
 D_0(x)\ge\delta
 \qquad (0\le t\le t_1,\ x\in\mathcal O(t)).
\end{equation}
\end{lemma}

\begin{proof}
The distance $D_0$ is $1$-Lipschitz. Fix $4<\beta'<\beta_0$ and a
positive absolute constant $c$, and consider
\[
 B(t,x,v)=\exp(C_1t)\exp\!\left(
       -(\beta_0-C_1t)|v|^2-\frac{cD_0(x)^2}{t}\right).
\]
At almost every point of a free path,
\begin{align*}
 \D\log B
 &=C_1+C_1|v|^2+\frac{cD_0^2}{t^2}
                  -\frac{2cD_0}{t}\,v\cdot\nabla D_0\\
 &\ge C_1+(C_1-c)|v|^2.
\end{align*}
The last inequality is completion of the square in $D_0/t$.
At fixed $(t,x)$, $B$ is a Maxwellian in velocity. Energy conservation in a
collision therefore gives $P(B,B)=B\nu_B$, and, as long as
$\beta_0-C_1t\ge\beta'>4$ and $e^{C_1t}\le2$,
\[
 P^b(B,B)\le P(B,B)\le C(1+|v|)B.
\]
Choose $C_1$ large enough that the preceding lower bound for
$\D\log B$ dominates $C(1+|v|)$, and then choose $t_1\le1/32$ small
enough to maintain both restrictions. These choices use only $\beta_0$
and the collision normalization.

The trace of $B$ along a path starting outside the hole dominates $u^0$.
Indeed, if $D_0(x)=0$, then $D_0(x+tv)\le t|v|$, so the distance factor
tends to one. If $D_0(x)>0$, its trace is zero, as is $u^0(x,v)$.
The differential inequality can be integrated from a positive time and
then passed to this trace. The positive-gain comparison used for
Proposition~\ref{cold:control}, with the nonnegative loss omitted, yields $u\le B$.
This uses only absolute continuity along paths at positive times and
boundedness in the $H$ norm; it does not differentiate the initial spatial
mask. Since $e^{C_1t}$ is bounded, \eqref{early:bath-bound} follows.

For $t\le1/32$, the axial interval defining $\mathcal O_i(t)$ has
half-length
\[
 a+2t+b_{\rm sp}\le\frac18+\frac1{16}+\frac1{16}=\frac14<\frac13.
\]
Its transverse distance from the corresponding lifted axis is at most
$s_0$. Thus every point of $\mathcal O(t)$ lies within distance $s_0$ of
the segment defining the hole. Since the hole is its open $2\delta$
neighborhood, the triangle inequality gives
$D_0\ge2\delta-s_0\ge\delta$ when $s_0\le\delta$.
\end{proof}

We fix $t_1$ from now on. Recall that
\[
 \mathcal S=2P^b(c,u)+\sum_{i\ne i'}P^b(c_i,c_{i'})
\]
is the source in the hot equation: it consists of collisions between
different cold colors and collisions between a cold color and the bath.

\begin{lemma}[Source estimates]\label{early:source-estimates}
For sufficiently small $s_0$ depending on $\delta$, every finite evolution
satisfies
\begin{equation}\label{early:source}
 P(c+u,H)\le\frac{C}{t+s_N}H,
 \qquad
 0\le\mathcal S\le C_{A,\delta}t^{A-1}H,
 \qquad 0<t\le t_1,
\end{equation}
on its interval of existence, for every fixed $A>0$.
Here $C$ is independent of $\delta$ and $R_\mu$, and the constants are
uniform for all sufficiently small $s_0$. No bound on $h$ is assumed.
\end{lemma}

\begin{proof}
At each $(t,x)$ there is at most one main color. Denote it by $a$, taking
$a=0$ if no main color is present, and put $r=c-a$.
The main-color measure bound in Lemma~\ref{cold:core-amplitude} and the sheet
gain estimates, including Lemma~\ref{ker:parameter-sheets}, give
\[
 P(a,H)\le\frac{C}{t+s_N}H.
\]
By \eqref{cold:nonmain-bound}, for each $M>0$,
\[
 0\le r\le\rho(t)H,\qquad \rho(t)\le C_Mt^M,
\]
uniformly for sufficiently small $s_0$. The bath bound gives $u\le CH$.
The quadratic gain bound of Lemma~\ref{lem:hot-bound} proves the first
estimate in \eqref{early:source}, absorbing the bounded terms since
$t+s_N\le2$.

If $a\ne0$, \eqref{early:hole-distance} improves the bath estimate to
$u\le C e^{-c\delta^2/t}H$ at this position. The ordered cross-color sum
is bounded by $2P(a,r)+P(r,r)$, so positivity yields
\[
 \mathcal S\le
 C\left(\frac{e^{-c\delta^2/t}+\rho(t)}{t+s_N}
                 +\rho(t)+\rho(t)^2\right)H.
\]
The exponential is bounded by any prescribed power of $t$, with a
constant depending on that power and on $\delta$. Taking $M=A$ proves
the second estimate. When $a=0$, only the last two types of terms remain,
and the same conclusion follows.
\end{proof}

Set $Z(t)=\|h(t)\|_H$. Dropping the loss in the hot path equation and
using \eqref{early:source}, we obtain, on every interval where $Z\le1$,
\begin{equation}\label{early:hot-comparison}
 Z(t)\le Z(r)+\int_r^t\left(
       \frac{C_2}{s+s_N}Z(s)+C_{A,\delta}s^{A-1}\right)\dd s,
 \qquad 0\le r\le t\le t_1.
\end{equation}
The constant $C_2$ includes the quadratic hot gain, because $Z^2\le Z$
and $s+s_N\le2$. It is independent of $\delta$, $R_\mu$, $s_0$, and $N$.
We will use integral inequalities such as \eqref{early:hot-comparison},
without assuming differentiability of the norm $Z$.

\subsection{Selecting a solution that touches the cap}

Fix the exponent in the seed \eqref{con:seed} so that
\begin{equation}\label{early:exponents}
 p>C_2+1,
 \qquad A>p+C_2+3,
 \qquad m_T=T^p.
\end{equation}
We call the solution with zero hot seed the \emph{dormant solution}.
The following selection adapts the cap argument of
\cite[Section~5]{OpenAIPeriodic}; we give the continuity argument here for
the present cutoff and color system.

\begin{proposition}[Cap selection]\label{early:cap-selection}
For fixed $\delta>0$, there are $\bar s_0,T_0>0$, uniform in $N$ and
the collision cutoff, with the following property.
For each $0<s_0\le\bar s_0$, each $0<T\le T_0$, and each finite
approximation, a seed parameter
$\lambda\in(0,L_0)$ gives a solution on $[0,T]$ satisfying
\begin{equation}\label{early:cap}
 Z(t)\le m_T(t+s_N)\quad(0\le t\le T),
 \qquad
 \max_{0\le t\le T}\frac{Z(t)}{m_T(t+s_N)}=1.
\end{equation}
The dormant solution exists on a fixed initial interval and satisfies
\begin{equation}\label{early:dormant-bound}
 Z(t)\le C'_{A,\delta}t^A.
\end{equation}
The interval and constant in \eqref{early:dormant-bound} are independent
of $N$, the collision cutoff, and sufficiently small $s_0$.
\end{proposition}

\begin{proof}
For the dormant solution, the function $Y(t)=C_dt^A$ is a supersolution
of \eqref{early:hot-comparison} if
\[
 C_d(A-C_2)\ge C_{A,\delta},
\]
because $(t+s_N)^{-1}\le t^{-1}$. At each fixed $N$, comparison begins
at the regular initial section $t=0$, where both functions vanish.
Choosing an initial interval on which $C_dt^A<1/2$ closes the condition
$Z\le1$. The cold and bath bounds are finite in the $H$ norm at each
fixed approximation. Thus the continuation criterion of
Proposition~\ref{con:finite-evolution} also proves existence on that
interval and gives \eqref{early:dormant-bound}.

Reduce $T_0$ until this interval contains $[0,T_0]$ and
$C_dT_0^{A-p-1}<1$. Then, for $0<t\le T\le T_0$,
\[
 \frac{Z(t)}{T^p(t+s_N)}
 \le\frac{C_dt^{A-1}}{T^p}
 \le C_dT^{A-p-1}<1.
\]
At $t=0$ the ratio is zero. Hence the dormant solution stays strictly
below the cap. At the other endpoint, $\lambda=L_0$, the ratio at time
zero equals $L_0\|\chi\|_H>1$.

Fix $T,N$, and the cutoff. Let $I$ be the set of
$\lambda\in[0,L_0]$ whose solutions exist on $[0,T]$ and remain strictly
below the cap there. Continuous dependence in the local theory and
continuation at the terminal section make $I$ relatively open.
It contains $0$ and does not contain $L_0$. Let $\lambda_*$ be the
right endpoint of the connected component of $I$ containing zero.
Along a sequence $\lambda_k\uparrow\lambda_*$ within that component,
all hot parts are bounded by the cap. The cold and bath estimates bound
the remaining colors uniformly in $k$ in the fixed-approximation norm.
The difference estimate in Proposition~\ref{con:finite-evolution}
therefore makes these solutions Cauchy, uniformly on $[0,T]$ in free
coordinates. Their limit is the solution with parameter $\lambda_*$
and satisfies the non-strict cap inequality.

The norm $Z$ is continuous in time, and the cap denominator is positive,
so the displayed ratio has a maximum. If that maximum were less than
one, continuous dependence and continuation would place a neighborhood
of $\lambda_*$ in $I$, contrary to its endpoint property. Its maximum
is therefore one. Since $0$ has a neighborhood in $I$, while
$L_0\|\chi\|_H>1$ also excludes a neighborhood of $L_0$, one has
$0<\lambda_*<L_0$. A final reduction of $T_0$ makes the cap at most one
on $[0,T]$ for every $T\le T_0$.
\end{proof}

\subsection{Smallness through the fixed bath time}

The pointwise estimate in \eqref{early:source} is very strong near zero,
but its constant depends on the hole width. To reach the fixed time
$t_1$ with an arbitrarily small hot part, we also use the shrinking
spatial support of the main colors.

\begin{lemma}[Small source along free paths]\label{early:path-source}
Fix $0<t_b<t_1$. There is a function $E_{t_b}(s_0)$ tending to zero as
$s_0\to0$ such that, for every finite evolution and every terminal time
$t\in[t_b,t_1]$ in its interval of existence,
\begin{equation}\label{early:source-integral}
 \left\|\int_{t_b}^t
    \mathcal S(s,x-(t-s)v,v)\dd s\right\|_H
 \le E_{t_b}(s_0).
\end{equation}
This estimate is uniform in the hot seed, the cap horizon, the
approximation index, and the collision cutoff.
\end{lemma}

\begin{proof}
On $[t_b,t_1]$, the nonmain bound \eqref{cold:nonmain-bound} gives
$r\le\varepsilon(s_0)H$, with $\varepsilon(s_0)\to0$.
The main-color sheet amplitude is bounded by $C/t_b$.
Consequently all terms of $\mathcal S$ containing a nonmain factor
are bounded by $C_{t_b}\varepsilon(s_0)H$, after enlarging
$\varepsilon$ if needed.

For the remaining main--bath term, define a deterministic majorant
\[
 \mathcal U_{\mathrm{main}}(s,x,v)
   =\sum_i\ind_{\mathcal O_i(s)}(x)\mathcal U_i(s,x,v),
 \qquad U_{\mathrm{bath}}(v)=C e^{-\beta'|v|^2},
\]
\[
 \mathcal S_{\mathrm{mb}}(s,x,v)
   =2P(\mathcal U_{\mathrm{main}}(s,x,\cdot),U_{\mathrm{bath}})(v).
\]
Here the sum includes all scales, so it bounds every finite selection;
at each position it has at most one nonzero summand.
By Proposition~\ref{cold:control} and \eqref{early:bath-bound},
\begin{equation}\label{early:common-majorant}
 0\le\mathcal S(s,x,v)
 \le C_{t_b}\varepsilon(s_0)H(v)+\mathcal S_{\mathrm{mb}}(s,x,v),
 \qquad t_b\le s\le t_1.
\end{equation}
The sheet gain estimates give $\mathcal S_{\mathrm{mb}}\le C_{t_b}H$.
The uniform large-output estimate \eqref{ker:large-output} gives
\begin{equation}\label{early:output-tail}
 \sup_{s,x,\,|v|>R}\frac{\mathcal S_{\mathrm{mb}}(s,x,v)}{H(v)}
 \le\omega_{t_b}(R),\qquad \omega_{t_b}(R)\longrightarrow0
 \quad(R\to\infty),
\end{equation}
with essential suprema in the velocity variable.
To handle velocities near the axis, fix $q'\in(1,q)$ and write
\[
 H_{q'}(v)=e^{-|v|^2}\langle v\rangle^{-100}
                 (1+|\bar v|^{-q'}).
\]
Since $U_{\mathrm{bath}}\le C_{q'}H_{q'}$, the same sheet lemma with exponent $q'$
implies $\mathcal S_{\mathrm{mb}}\le C_{t_b,q'}H_{q'}$. Hence, for $0<r\le1$,
\begin{equation}\label{early:output-axis}
 \sup_{s,x,\,0<|\bar v|<r}\frac{\mathcal S_{\mathrm{mb}}(s,x,v)}{H(v)}
 \le C_{t_b,q'}r^{q-q'}.
\end{equation}
Both estimates include the Gaussian-parameter measures in
Lemma~\ref{cold:core-amplitude}, by Lemma~\ref{ker:parameter-sheets}.

It remains to consider the compact velocity region
$|v|\le R$, $|\bar v|\ge r$. The spatial support of $\mathcal S_{\mathrm{mb}}$ is contained,
for every $s\le t_1$, in the projection of the fixed union
\[
 \bigcup_{n\in\Lambda}
     \bigl(n+B^2_{s_0}\times[-1/4,1/4]\bigr),
\]
where $B^2_{s_0}$ is the transverse disk of radius $s_0$.
A lifted free path with $|v|\le R$ and duration at most $t_1$ can meet
only $C_R$ of these cylinders, uniformly in its endpoint: their centers
must lie in a fixed bounded enlargement of a segment of length at most
$Rt_1$, and the lattice has a fixed separation.
For any one cylinder, its transverse projection is crossed in time at
most $2s_0/|\bar v|\le2s_0/r$.
Thus the normalized integral of $\mathcal S_{\mathrm{mb}}$ along such a path is bounded by
$C_{t_b,R}s_0/r$. Combining the three velocity regions gives
\[
 \left\|\int_{t_b}^t
     \mathcal S(s,x-(t-s)v,v)\dd s\right\|_H
 \le C_{t_b}\varepsilon(s_0)
      +t_1\omega_{t_b}(R)
      +C_{t_b,q'}t_1r^{q-q'}
      +C_{t_b,R}\frac{s_0}{r}.
\]
First take $R$ large, then $r$ small, and finally $s_0$ small.
This proves \eqref{early:source-integral}.

All estimates above concern the common pointwise majorant
\eqref{early:common-majorant}, rather than values chosen separately on
individual paths. To pass an almost-everywhere inequality to path
integrals, use free coordinates $(s,x_0,v)\mapsto(s,x_0+sv,v)$.
This map preserves measure, so Fubini's theorem gives the inequality at
almost every time on almost every free path. Integration on those paths
then defines representatives at every terminal section. Since free
transport preserves $\|\cdot\|_H$, the resulting bound is precisely the
essential norm in \eqref{early:source-integral} at each fixed terminal
time. The construction applies on any existing solution interval, so
this lemma does not presuppose continuation to $t_1$.
\end{proof}

\begin{proposition}[Uniform early continuation]\label{early:continuation}
Fix $\delta>0$ and $\zeta\in(0,1)$.
There are $\bar s_0>0$ and $T_0\in(0,t_1)$, depending on $\delta$ and
$\zeta$ and on the fixed jet parameters, such that for every
$0<s_0\le\bar s_0$, every finite approximation, and every
$0<T\le T_0$, the dormant solution and every solution selected in
Proposition~\ref{early:cap-selection} continue to $t_1$ and satisfy
\begin{equation}\label{early:small-hot}
 \sup_{0\le t\le t_1}\|h(t)\|_H\le\zeta.
\end{equation}
The choices are uniform in $N$, the collision cutoff, $R_\mu$, and $T$.
For dormant solutions, \eqref{early:dormant-bound} also holds on an
initial interval independent of these indices, with the fixed exponent
$A>p+1$ chosen in \eqref{early:exponents}.
\end{proposition}

\begin{proof}
We choose an intermediate time $t_b\in(0,t_1)$ and then require
$T_0<t_b$. At the cap horizon, both the selected and dormant solutions
satisfy
\[
 Z(T)\le T^p(T+s_N)\le2T^p,
\]
using only $T,s_N\le1$; no relation between $T$ and $s_N$ is needed.
For $T\le t\le t_b$, \eqref{early:hot-comparison} and
$(s+s_N)^{-1}\le s^{-1}$ give, while $Z\le1$,
\begin{equation}\label{early:intermediate-bound}
 Z(t)\le2T^{p-C_2}t^{C_2}
                  +\widetilde C_{A,\delta}t^A.
\end{equation}
Indeed multiplication of the scalar comparison by $t^{-C_2}$ integrates
the source to at most
$C_{A,\delta}t^{A-C_2}/(A-C_2)$.

From $t_b$ onward, retain the source inside the path integral rather
than using its pointwise bound. The same hot gain estimates and
Lemma~\ref{early:path-source} give
\[
 Z(t)\le Z(t_b)+E_{t_b}(s_0)
                 +\int_{t_b}^t\frac{C_2}{s}Z(s)\dd s.
\]
Consequently, as long as $Z\le1$,
\begin{equation}\label{early:final-comparison}
 \sup_{t_b\le t\le t_1}Z(t)
 \le\left(\frac{t_1}{t_b}\right)^{C_2}
       \left(2T_0^{p-C_2}t_b^{C_2}
              +\widetilde C_{A,\delta}t_b^A+E_{t_b}(s_0)\right),
\end{equation}
where the supremum is restricted to the current interval of existence
until continuation has been established.

The order of choices now matters. Since $A>C_2$, choose $t_b$ so that
\[
 \widetilde C_{A,\delta}t_1^{C_2}t_b^{A-C_2}<\frac\zeta6.
\]
Since $p>C_2$, reduce $T_0<t_b$, within the range of
Proposition~\ref{early:cap-selection}, so that
\[
 2T_0^{p-C_2}t_1^{C_2}<\frac\zeta6,
 \qquad 2T_0^p<\frac\zeta2.
\]
Finally take $\bar s_0$ small enough for all preceding cold and bath
estimates and for
\[
 \left(\frac{t_1}{t_b}\right)^{C_2}E_{t_b}(s_0)<\frac\zeta6
 \qquad(0<s_0\le\bar s_0).
\]
The last choice is possible by Lemma~\ref{early:path-source}.
Equation~\eqref{early:final-comparison} is then strictly below
$\zeta/2$. The same first two choices make
\eqref{early:intermediate-bound} strictly below $\zeta/2$ on
$[T,t_b]$, and the cap gives the same bound on $[0,T]$.

These estimates close the condition $Z\le1$ with a strict margin.
At each fixed approximation the cold sum and bath have bounded
$H$ norm through $t_1$, so the finite-evolution continuation criterion
excludes a finite endpoint before $t_1$.
This proves existence and \eqref{early:small-hot} uniformly in all the
stated indices. The dormant estimate was obtained independently on its
initial interval and is preserved under this continuation.
\end{proof}
